\begin{lemma}
\label{lemma-etale-makes-quasi-finite-finite-one-prime}
Let $R \to S$ be a ring map.
Let $\mathfrak q \subset S$ be a prime lying over
the prime $\mathfrak p \subset R$.
Assume $R \to S$ finite type and quasi-finite at $\mathfrak q$.
Then there exists
\begin{enumerate}
\item an \'etale ring map $R \to R'$,
\item a prime $\mathfrak p' \subset R'$ lying over $\mathfrak p$,
\item a product decomposition
$$
R' \otimes_R S = A \times B
$$
\end{enumerate}
with the following properties
\begin{enumerate}
\item $\kappa(\mathfrak p) = \kappa(\mathfrak p')$,
\item $R' \to A$ is finite,
\item $A$ has exactly one prime $\mathfrak r$ lying over $\mathfrak p'$,
\item $\mathfrak r$ lies over $\mathfrak q$, and
\item $B$ does not have a prime lying over $\mathfrak q$ and $\mathfrak p'$.
\end{enumerate}
\end{lemma}

\begin{proof}
Let $S' \subset S$ be the integral closure of $R$ in $S$.
Let $\mathfrak q' = S' \cap \mathfrak q$.
By Zariski's Main Theorem \ref{theorem-main-theorem}
there exists a $g \in S'$, $g \not \in \mathfrak q'$ such
that $S'_g \cong S_g$. Consider the fibre rings
$F = S \otimes_R \kappa(\mathfrak p)$ and
$F' = S' \otimes_R \kappa(\mathfrak p)$. Denote $\overline{\mathfrak q}'$
the prime of $F'$ corresponding to $\mathfrak q'$. Since
$F'$ is integral over $\kappa(\mathfrak p)$ we see
that $\overline{\mathfrak q}'$ is a closed point of
$\Spec(F')$, see Lemma \ref{lemma-integral-over-field}.
Note that $\mathfrak q$ defines an isolated closed point
$\overline{\mathfrak q}$ of
$\Spec(F)$ (see Definition \ref{definition-quasi-finite}).
Since $S'_g \cong S_g$ we have $F'_g \cong F_g$,
so $\overline{\mathfrak q}$ and $\overline{\mathfrak q}'$
have isomorphic open neighbourhoods in $\Spec(F)$
and $\Spec(F')$. We conclude the set
$\{\overline{\mathfrak q}'\} \subset \Spec(F')$ is
open. Combined with $\mathfrak q'$ being closed (shown above)
we conclude that $\overline{\mathfrak q}'$ defines
an isolated closed point of $\Spec(F')$ as well.

\medskip\noindent
An additional small remark is that under the map
$\Spec(F) \to \Spec(F')$ the point $\overline{\mathfrak q}$
is the only point mapping to $\overline{\mathfrak q}'$. This follows
from the discussion above.

\medskip\noindent
By Lemma \ref{lemma-disjoint-implies-product} we may write
$F' = F'_1 \times F'_2$ with
$\Spec(F'_1) = \{\overline{\mathfrak q}'\}$.
Since $F' = S' \otimes_R \kappa(\mathfrak p)$, there
exists an $s' \in S'$ which maps to the element
$(r, 0) \in F'_1 \times F'_2 = F'$ for some $r \in R$, $r \not \in \mathfrak p$.
In fact, what we will use about $s'$ is that it is an element of $S'$,
not contained in $\mathfrak q'$, and contained in any other prime
lying over $\mathfrak p$.

\medskip\noindent
Let $f(x) \in R[x]$ be a monic polynomial such that $f(s') = 0$.
Denote $\overline{f} \in \kappa(\mathfrak p)[x]$ the image.
We can factor it as $\overline{f} = x^e \overline{h}$ where
$\overline{h}(0) \not = 0$. After replacing $f$
by $x f$ if necessary, we may assume $e \geq 1$.
By Lemma \ref{lemma-factor-mod-lift-etale}
we can find an \'etale ring extension $R \to R'$,
a prime $\mathfrak p'$ lying over $\mathfrak p$, and
a factorization $f = h i$ in $R'[x]$ such that
$\kappa(\mathfrak p) = \kappa(\mathfrak p')$,
$\overline{h} = h \bmod \mathfrak p'$,
$x^e = i \bmod \mathfrak p'$, and
we can write $a h + b i = 1$ in $R'[x]$ (for suitable $a, b$).

\medskip\noindent
Consider the elements $h(s'), i(s') \in R' \otimes_R S'$.
By construction we have $h(s')i(s') = f(s') = 0$. On the other
hand they generate the unit ideal since $a(s')h(s') + b(s')i(s') = 1$.
Thus we see that $R' \otimes_R S'$ is the product of the
localizations at these elements:
$$
R' \otimes_R S'
=
(R' \otimes_R S')_{i(s')}
\times
(R' \otimes_R S')_{h(s')}
=
S'_1 \times S'_2
$$
Moreover this product decomposition is compatible with the product
decomposition we found for the fibre ring $F'$; this comes from our
choices of $s', i, h$ which guarantee that $\overline{\mathfrak q}'$
is the only prime of $F'$ which does not contain the image of $i(s')$
in $F'$. Here we use that the fibre ring of $R'\otimes_R S'$ over $R'$ at
$\mathfrak p'$ is the same as $F'$ due to the fact that
$\kappa(\mathfrak p) = \kappa(\mathfrak p')$.
It follows that $S'_1$  has exactly
one prime, say $\mathfrak r'$,
lying over $\mathfrak p'$ and
that this prime lies over $\mathfrak q'$.
Hence the element $g \in S'$ maps to an element of $S'_1$ not contained
in $\mathfrak r'$.

\medskip\noindent
The base change $R'\otimes_R S$ inherits a similar product decomposition
$$
R' \otimes_R S
=
(R' \otimes_R S)_{i(s')}
\times
(R' \otimes_R S)_{h(s')}
=
S_1 \times S_2
$$
It follows from the above that $S_1$ has exactly
one prime, say $\mathfrak r$,
lying over $\mathfrak p'$ (consider the fibre ring as above),
and that this prime lies over $\mathfrak q$.

\medskip\noindent
Now we may apply Lemma \ref{lemma-produce-finite} to the ring maps
$R' \to S'_1 \to S_1$, the prime $\mathfrak p'$ and
the element $g$ to see that after replacing $R'$ by
a principal localization we can assume that $S_1$ is
finite over $R'$ as desired.
\end{proof}
